% !TeX root = ../main.tex

\chapter{Architettura del Sistema e Osservabilità}
\label{cap:architettura_osservabilita}

Nel progetto l'architettura è basata su \texttt{Kind} (Kubernetes IN Docker), viene utilizzato per avviare cluster Kubernetes locali e leggeri dentro dei container di Docker replicando molti comportamenti di un cluster reale (scheduler, networking, controller) in un ambiente contenuto e facilmente ricreabile, ideale per validazione e test prima di un eventuale rilascio.

Kubernetes è un orchestratore di container che gestisce deployment, scaling e resilienza delle applicazioni: organizza i container in Pod, mantiene lo stato desiderato con Deployment e ReplicaSet, espone servizi con Service e fornisce autoscaling e self-healing.
Aspetti che nel progetto vengono utilizzati tramite invocazione delle funzioni di \texttt{kubectl}.

Docker è una piattaforma di containerizzazione che confeziona applicazioni e dipendenze in immagini portabili, garantisce isolamento, avvii rapidi e coerenza tra ambienti.
Funge da base su cui costruire e distribuire le immagini eseguite da kind (o Kubernetes).

\section{Design a Microservizi}
\label{sec:design_microservizi}
Nel progetto abbiamo adottato un design a microservizi che separa chiaramente responsabilità e domini di failure in tre blocchi indipendenti: Data Layer (InfluxDB) dedicato alla persistenza di metriche/serie temporali, Message Broker (MQTT/Mosquitto) come canale asincrono di scambio eventi, e Logic Layer (Server Centrale e servizi applicativi) che implementa l’elaborazione e l’esposizione delle funzionalità.

Questa separazione è anche fisica perché ogni blocco viene eseguito in Pod distinti (con Service per discovery via DNS interno, es. \texttt{influxdb-service} e \texttt{mosquitto-service}), così che il layer logico non dipenda da IP/staticità ma solo da endpoint di servizio.

Il broker disaccoppia produttori e consumatori: il layer logico può scalare o riavviarsi senza “spezzare” l’ingest, mentre InfluxDB resta un componente stateful con volume persistente (PVC) e strategia di avvio che privilegia consistenza dei dati.

Gli “script” Python non vengono eseguiti come processi manuali sui nodi: Kubernetes li avvia come processo principale del container tramite il command del Deployment (un container = un servizio, oppure più container nello stesso Pod quando serve co-locazione, ad esempio per condividere dati/volumi e ridurre latenza).

Lo scheduling decide su quale nodo far girare ciascun Pod; la replica (campo \texttt{replicas}) permette di distribuire più istanze di uno stesso microservizio su nodi diversi per aumentare throughput e tolleranza ai guasti, mantenendo invariata l’interfaccia grazie ai Service.

\section{Simulazione del Mondo Fisico (Digital Twin)}
\label{sec:digital_twin}

Nel progetto i digital twin di droni e client sono implementati come generatori controllati di telemetria e domanda, utili a validare l’architettura end‑to‑end senza dipendere da dispositivi fisici.

Il twin del drone (in \texttt{drone\_simulator.py}) mantiene uno stato interno minimale ma realistico (posizione, batteria, usura e stato operativo) che evolve a tick: pubblica periodicamente la telemetria sul broker MQTT (topic \texttt{telemetry/drones}) e si mette in ascolto di comandi sul proprio canale dedicato comandi \texttt{DRONE\_ID}, reagendo a events come assegnazione missione o rientro alla base e simulando consumi/degrado fino a condizioni di \texttt{MAINTENANCE}.

In parallelo, il twin del client (in \texttt{client\_simulator.py}) genera carico “business” pubblicando nuovi ordini su MQTT (topic \texttt{business/ordini/nuovi}) con intervalli casuali e payload strutturati (coordinate pickup/drop, peso, priorità), così da stressare broker e layer applicativo in modo riproducibile.

Dal punto di vista infrastrutturale, entrambi vengono eseguiti su Kubernetes como Deployment separati, ciascuno nel proprio Pod (configurati nei manifest \texttt{configs/drone.yaml} e \texttt{configs/client.yaml}): il processo principale del container è lo script Python definito nel command, e la connettività verso l’infrastruttura è ottenuta via service discovery (\texttt{mosquitto-service}).

La “moltiplicazione” della flotta è ottenuta in modo nativo tramite lo scaling: aumentando \texttt{replicas} del Deployment drone-simulator Kubernetes crea più Pod identici, e poiché \texttt{DRONE\_ID} viene derivato dal \texttt{metadata.name} del Pod, ogni replica corrisponde a un nuovo drone digitale con identità e canale comandi distinti.

Questo scaling può essere attivato anche in modo automatico dall’agente: attraverso il server MCP viene esposto il tool '\texttt{scale\_drone\_deployment}' che invoca la patch della risorsa '\texttt{deployment/scale}' (equivalente concettuale di un '\texttt{kubectl scale}'), con guardrail operativi per mantenere controllato l'impatto delle azioni automatiche sul cluster 'vedi capitolo 4'.

\section{Osservabilità (Observability Strategy)}
\label{sec:osservabilita}
Come da direttive del corso, un sistema distribuito deve permettere una chiara visibilità operativa per consentire il monitoraggio dello stato, la tracciabilità delle decisioni agentiche e la rapida mitigazione dei guasti. Nel nostro progetto una strategia di osservabilità basata su tre pilastri fondamentali: la visualizzazione real-time tramite una Dashboard Web centralizzata, il tracciamento di log strutturati per la tracciabilità delle azioni e la persistenza di metriche temporali.

\subsection*{1. Dashboard Web Centralizzata}
La visibilità generale dello stato del sistema è affidata al server centrale (\texttt{central\_server.py}), il quale ospita un'interfaccia utente web integrata tramite template HTML asincroni e un set di endpoint REST nativi in Flask: 

\begin{itemize}
    \item \texttt{/api/status}: Fornisce un riepilogo numerico aggregato (\textit{Snapshot}) del sistema, permettendo la visualizzazione del conteggio totale dei droni attivi, degli ordini pendenti e delle consegne completate, oltre alla configurazione dei nodi di rete (broker MQTT e istanza InfluxDB).
    \item \texttt{/api/orders}: Espone la coda dei messaggi in stato \texttt{PENDING} memorizzati in RAM.
    \item \texttt{/api/drones}: Restituisce la matrice completa dei droni agganciati al cluster, visualizzandone dinamicamente lo stato e le altre varie metriche.
    \item \texttt{/api/completed}: Elenca lo storico logistico delle consegne andate a buon fine.
\end{itemize}

Lato client, la dashboard esegue un ciclo di \textit{polling} asincrono tramite chiamate \texttt{fetch()} JavaScript ogni 5000 millisecondi (5 secondi). Questo intervallo garantisce un aggiornamento visivo costante.

\subsection*{2. Log di Sistema e Audit Trail dei Comandi Agentici}
Il tracciamento dei log è strutturato su più livelli per garantire l'ispezionabilità (\textit{Auditability}) richiesta.

\begin{enumerate}
    \item \textbf{Registro di Audit Strutturato (\texttt{audit\_actions.jsonl}):} Gestito nativamente dal modulo \texttt{drone\_mcp\_layer.py}, archivia in formato JSON Lines ogni singola operazione di scrittura o rimediazione avviata dall'IA. Traccia timestamp, il tipo di azione chiamata e l'esatto payload.
    \item \textbf{Registro delle Approvazioni Umane (\texttt{pending\_approvals.jsonl}):} Traccia il ciclo di vita delle richieste degli agenti quando riconoscono il superamento dei confini di policy. Il file contiene la richiesta iniziale dell'agente, la giustificazione logica fornita dall'LLM e il successivo cambiamento di stato causato dall'operatore.
    \item \textbf{Log di Runtime e Tracciamento dei Costi:} I simulatori dei droni stampano a terminale log continui riguardanti il loro stato. Parallelamente, l'orchestratore e gli agenti loggano per ciascuna iterazione l'esatto consumo di token restituito dai modelli. Questa metrica loggata è cruciale per la verifica dei tetti di spesa.
\end{enumerate}

\subsection*{3. Metriche di Serie Temporali su InfluxDB}
L'ultimo pilastro dell'osservabilità riguarda l'archiviazione e la persistenza dei dati operativi sotto forma di serie temporali. 
Quando il server centrale intercetta i messaggi sui topic MQTT, scrive in modalità sincrona record all'interno del bucket \texttt{iot\_data}:

\begin{itemize}
    \item \textbf{Measurement \texttt{drone\_telemetry}:} Registra l'evoluzione fisica della flotta.
    \item \textbf{Measurement \texttt{business\_orders}:} Monitora il flusso di ordini in ingresso.
\end{itemize}

Tramite il layer MCP (\texttt{get\_drones\_telemetry} e \texttt{get\_pending\_orders}), gli LLM possono analizzare finestre temporali passate (\textit{time-windows} di 5 o 60 minuti) per estrarre trend di carico o pattern di deterioramento hardware.